\documentclass[11pt]{article}
\usepackage{setspace}
\usepackage{geometry}
\geometry{margin=1in}
\setstretch{1.15}

\begin{document}

% --- Applicant Information ---
\noindent
\textbf{Zhongli Wang} \\
Phone: +49 160 9541 6015 \\
Email: nkwangzhongli@gmail.com \\
Website: https://sites.google.com/view/zhongli-wang \\
Bielefeld, Germany \\

\vspace{0.5cm}

% --- Recipient Information ---
\noindent
\textbf{Recruitment Committee} \\
Department of Economics \\
University of Vienna \\
Vienna, Austria \\

\vspace{0.5cm}

% --- Subject Line ---
\noindent
\textbf{Application for University Assistant (Postdoctoral) — Applied Microeconomics (Job ID: 4860)} \\

\vspace{0.3cm}

\noindent Dear Prof. Pichler and Members of the Recruitment Committee,\\[0.2cm]

I am writing to apply for the position of University Assistant (postdoctoral) in Applied Microeconomics at the University of Vienna (Job ID: 4860). I am completing my Ph.D.\ in Economics at Bielefeld University and Università Cattolica del Sacro Cuore through the Marie Skłodowska–Curie EPOC programme. My research focuses on applied microeconomics with a specialization in digital markets, industrial organization, and computational modeling. I combine structural modeling, simulation-based econometrics, and machine-learning tools to study firm behaviour, market design, and the interaction between technological change and economic decision-making. This position is an excellent match for my research profile and long-term academic goals.

My dissertation, \textit{“Strategic Behavior and Economic Dynamics in ICT Markets: An Agent-Based Framework,”} develops a series of applied microeconomic models investigating how firms compete in digital environments characterized by network effects, congestion, investment frictions, and heterogeneous consumers. I formalize firm decisions in pricing, capacity investment, and quality provision, and link them to market outcomes through empirically validated behavioural and structural parameters. To estimate these models, I employ the Simulated Method of Moments, Monte Carlo simulations, and neural-network surrogate models. These approaches allow me to analyse counterfactual policy scenarios and derive comparative-statics insights for markets undergoing rapid technological transformation. The project contributes to digital economics and empirical IO by offering transparent modeling tools that integrate firm heterogeneity, dynamic competition, and computational estimation.

A unifying theme of my work is understanding how micro-level behavioural mechanisms translate into aggregate outcomes—whether in market concentration, investment incentives, or welfare distribution. These questions align strongly with the Department of Economics’ strengths in applied micro, industrial organization, and macro-micro modeling. My work complements research on dynamic decision-making under uncertainty, bank behaviour and investment risk (e.g., Prof.\ Paul Pichler’s work), and market design under frictions. I am eager to deepen these connections and to expand my research into related areas such as digital platforms, energy markets, and the economics of technological regulation.

In addition to my dissertation, I have coauthored projects demonstrating my ability to work with large-scale data, high-dimensional models, and advanced quantitative tools. \textit{Polytope Fraud Theory} employs machine learning to identify abnormal financial patterns across firms, illustrating my experience with predictive modeling and empirical analysis. A second project develops AI-assisted research workflows using large-language-model agents, showcasing my commitment to reproducible, transparent, and computationally rigorous research methods. These skills enhance my ability to contribute to collaborative research projects within the department.

I also bring experience relevant to teaching and academic service. I have delivered invited talks on AI and economics (including in German), organized academic workshops, supervised student assistants, and mentored peers on computational techniques. I am prepared to independently teach courses in applied microeconomics, industrial organization, quantitative methods, and computational economics, and I am highly motivated to contribute to the department’s teaching mission. I am currently improving my German and am committed to further strengthening my language skills.

The University of Vienna’s internationally recognized research environment, its strong emphasis on applied microeconomics, and its supportive long-term postdoctoral structure make it an ideal setting to develop my independent research agenda. 
I am particularly motivated by the opportunity to collaborate with faculty whose work intersects IO, dynamic modeling, and the economics of regulation, and to contribute to a vibrant scholarly community.

Thank you very much for considering my application. I would be honored to join the Department of Economics at the University of Vienna and would welcome the opportunity to discuss my candidacy in more detail.

\vspace{6mm}

\noindent Sincerely,\\[12mm]
\noindent \textbf{Zhongli Wang}

\end{document}
